\section{Data Description}
\label{sec:data}

\subsection{Base Dataset: ADE20K}

ADE20K~\cite{zhou2017scene} is a large-scale scene parsing benchmark
providing over 20,000 images with dense 150-class pixel-level
annotations across diverse indoor and outdoor scenes. We use
3,000~images from the combined official training and validation
splits, selected for indoor scene coverage across kitchens, bathrooms,
living rooms, corridors, offices, retail spaces, and healthcare
environments.

ADE20K is selected as the base dataset for three reasons.
First, it provides broad indoor scene diversity with pixel-level
annotation, covering the environmental contexts most relevant to
domestic hazard understanding. Second, its 150-class annotation
scheme includes the floor, furniture, and object classes needed
for our hazard taxonomy remapping (Table~\ref{tab:taxonomy}).
Third, its widespread adoption enables comparison of our hazard
taxonomy with prior indoor segmentation work.

\subsection{ADE20K Hazard Remapping}

ADE20K's 150~semantic classes are remapped to our four-class hazard
taxonomy using fixed index mappings. Safe~Floor (class~1) includes
ADE20K indices 3 (floor), 6 (rug/carpet), 9 (staircase), and 11
(path). Obstacle (class~2) includes furniture indices
$\{5, 10, 14, 18, 22, 24, 30\}$ (chairs, tables, sofas, beds,
shelves, cabinets, desks) and clutter indices
$\{33, 34, 36, 39, 41, 43, 44, 62, 71, 88, 125\}$ (bags, boxes,
bins, suitcases, carts, and related objects).
Sharp~Objects (class~3) includes indices 96~(knife) and 97
(scissors). All remaining ADE20K classes map to Background (class~0).

\subsection{Manual Hazard Annotation Contribution}

Since no existing dataset provides pixel-level annotations aligned
with indoor injury risk, we introduce a \textbf{manual hazard
annotation dataset} of 100~ADE20K indoor images with pixel-level
labels created using Make~Sense~AI, a browser-based polygon
annotation tool (\url{makesense.ai}). Images were selected from
ADE20K to maximise hazard content: kitchen images for knife/scissors
presence, bathroom images for wet floor surfaces, and
living/dining room images for furniture and clutter obstacles.

\noindent\textbf{Annotation protocol.}
Three hazard classes were annotated using closed polygon regions:
\begin{itemize}
\item \textbf{Safe Floor} (mapped to class~1): All floor surfaces
including dry tiles, hardwood, carpets, and rugs.
\item \textbf{Obstacle} (mapped to class~2): Any furniture or loose
object that could cause a collision or trip hazard, including chairs,
tables, bags, and clutter on the floor.
\item \textbf{Sharp Objects} (mapped to class~3): Any visibly sharp
instrument including knives, scissors, and broken glass.
\end{itemize}

Table~\ref{tab:ann_stats} summarises annotation statistics.
The high number of Obstacle annotations (222) compared to Safe~Floor
(107) reflects the greater diversity of obstacle types in indoor
scenes. Sharp~Objects (186~annotations) is intentionally
overrepresented to address the extreme pixel-level sparsity of this
class in the ADE20K base annotations.

\begin{table}[h]
\centering
\caption{Manual Annotation Statistics (100 images, 199 annotated masks)}
\label{tab:ann_stats}
\begin{tabular}{lcc}
\toprule
\textbf{Class} & \textbf{Annotated Regions} & \textbf{Images Containing Class} \\
\midrule
Safe Floor & 107 & 72 \\
Obstacle & 222 & 89 \\
Sharp Objects & 186 & 63 \\
\midrule
\textbf{Total} & \textbf{515} &: \\
\bottomrule
\end{tabular}
\end{table}

\noindent\textbf{Label mapping.}
Exact string matching (case-insensitive) is applied to annotation
labels: ``sharp objects'' / ``sharp object'' / ``knife'' /
``scissors'' $\to$~3; ``obstacle'' / ``clutter'' /
``clutter/~obstacles'' / ``furniture'' $\to$~2;
``dry floor'' / ``wet floor'' / ``safe floor'' / ``floor'' $\to$~1.

\subsection{COCO Sharp Object Augmentation}

To address the extreme scarcity of Sharp~Objects pixels across
ADE20K, we incorporate 186~images from the COCO 2017 validation
set~\cite{lin2014coco} containing knife (category~49) or scissors
(category~87) instances with pixel-level instance segmentation
annotations. COCO instance masks are converted to class~3
(Sharp~Objects) semantic labels.

A key design decision is to apply \textbf{no floor-level spatial
constraint}. Prior work on indoor hazard detection has sometimes
argued that sharp objects are only hazardous when on the floor.
We reject this constraint: a knife on a kitchen counter can fall
and cause injury; scissors on a desk are accessible to children;
a sharp kitchen tool visible anywhere in the scene represents a
hazard worth flagging. Restricting labels to the lower portion
of the image would suppress valid positive annotations and
reduce training signal. We include all COCO sharp object
annotations regardless of vertical position, with a minimum
area threshold of 200~pixels to exclude degenerate instances.

\subsection{Mask Priority Merging}

Manual annotations override ADE20K remapping at annotated locations,
since manual labels are higher-quality and specifically designed for
hazard classification. COCO masks are used directly without any
ADE20K remapping. The priority order is:
(1) Manual annotation (highest priority),
(2) COCO annotation, (3) ADE20K remapping (fallback).

\subsection{Data Partitioning}

All 3,000~ADE20K images are shuffled with fixed seed~42 and
partitioned into a strict 70/15/15 split: 2,100~training,
450~validation, and 450~test images. Manual annotation images
and COCO images are added exclusively to the training set.
The 450~test images are used only for final evaluation and
are never used during model development or hyperparameter selection.

\subsection{Preprocessing and Augmentation}

All images and masks are resized to $256{\times}256$ pixels.
Masks use nearest-neighbour interpolation to prevent label
interpolation artefacts at class boundaries. Images are normalised
with ImageNet statistics (mean~$[0.485, 0.456, 0.406]$,
std~$[0.229, 0.224, 0.225]$). During training, random horizontal
flipping (probability~0.5) is applied with synchronised image-mask
transformation, doubling training sample diversity.

\subsection{Class Distribution and Imbalance Analysis}

Table~\ref{tab:class_dist} shows the approximate class distribution
in the training set. Background dominates at ${\sim}74\%$ of pixels,
reflecting the large proportion of walls, ceilings, and
unclassified surfaces in indoor scenes. Sharp~Objects represents
less than 1\% of total pixels despite 230~training images
containing the class (44~from ADE20K, 186~from COCO), an effective
imbalance ratio of approximately 46:1.

\begin{table}[h]
\centering
\caption{Training Set Class Distribution and Loss Weights}
\label{tab:class_dist}
\begin{tabular}{lcccc}
\toprule
\textbf{Class} & \textbf{Pixel \%} & \textbf{Weight} & \textbf{Images} & \textbf{Source} \\
\midrule
Background & $\sim$74\% & 0.5 & 2,100 & ADE20K \\
Safe Floor & $\sim$16\% & 1.5 & 1,847 & ADE20K+Manual \\
Obstacle & $\sim$9\% & 2.0 & 1,203 & ADE20K+Manual \\
Sharp Objects & $<1\%$ & 6.0 & 230 & ADE20K+Manual+COCO \\
\bottomrule
\end{tabular}
\end{table}

This extreme imbalance motivates the class weight of 6.0 for Sharp
Objects--12$\times$ the background weight--and the COCO augmentation
strategy. Despite these measures, Sharp~Objects achieves zero IoU
in evaluation, highlighting that pixel-level sparsity at this
scale ($<0.05\%$ of test pixels) represents a fundamental challenge
for discriminative training that class weighting alone cannot overcome.
